
\documentclass{article}
\usepackage{colortbl}
\usepackage{makecell}
\usepackage{multirow}
\usepackage{supertabular}

\begin{document}

\newcounter{utterance}

\centering \large Interaction Transcript for game `wordle', experiment `high\_frequency\_words\_no\_clue\_no\_critic', episode 2 with reflexion{-}Qwen3.5{-}9B.
\vspace{24pt}

{ \footnotesize  \setcounter{utterance}{1}
\setlength{\tabcolsep}{0pt}
\begin{supertabular}{c@{$\;$}|p{.15\linewidth}@{}p{.15\linewidth}p{.15\linewidth}p{.15\linewidth}p{.15\linewidth}p{.15\linewidth}}
    \# & \multicolumn{2}{c}{Player} && \multicolumn{2}{c}{Game Master} \\
    \hline

    \theutterance \stepcounter{utterance}  
    & & & \multicolumn{4}{p{0.6\linewidth}}{
        \cellcolor[rgb]{0.9,0.9,0.9}{
            \makecell[{{p{\linewidth}}}]{
                \texttt{\tiny{[P1$\langle$GM]}}
                \texttt{You are a language wizard who likes to guess words by using the given rules.} \\
\\ 
\texttt{Welcome to Wordle! You have six attempts to guess the target word, a valid English word of five lowercase letters (a{-}z). Please use the tags "explanation:" and "guess:" to provide a concise explanation for each guess.} \\
\\ 
\texttt{For instance, if your guess is "apple", your response should be} \\
\texttt{explanation: this is a common five{-}letter English word, and I am starting my guess with this word.} \\
\texttt{guess: apple} \\
\\ 
\texttt{After each guess, your answer will be validated, and you will receive feedback indicating which letters are correct (green), which letters are correct but in the wrong position (yellow), and which letters are incorrect (red). This feedback can be useful in determining which letters to include or exclude in your next guess.} \\
\\ 
\texttt{For example, the feedback for "apple" might be:} \\
\texttt{guess\_feedback: a<yellow> p<yellow> p<green> l<yellow> e<red>} \\
\\ 
\texttt{The explanation should contain details about how the guess\_feedback is used to arrive at a new guess.} \\
\\ 
\texttt{Do not generate any other text.} \\
\\ 
\texttt{Let's begin with your first guess.} \\
            }
        }
    }
     \\ \\

    \theutterance \stepcounter{utterance}  
    & \multicolumn{4}{p{0.6\linewidth}}{
        \cellcolor[rgb]{0.561,0.737,0.561}{
            \makecell[{{p{\linewidth}}}]{
                \texttt{\tiny{[P1$\rangle$GM]}}
                \texttt{explanation: I am starting with "crane" because it contains common vowels (a, e) and frequent consonants (r, n) that appear in many five{-}letter words, allowing me to quickly eliminate or confirm several letters.} \\
\texttt{guess: crane} \\
            }
        }
    }
    & & \\ \\

    \theutterance \stepcounter{utterance}  
    & & & \multicolumn{2}{p{0.3\linewidth}}{
        \cellcolor[rgb]{0.9,0.9,0.9}{
            \makecell[{{p{\linewidth}}}]{
                \texttt{\tiny{[GM$|$GM]}}
                \texttt{attempts = 1} \\
\texttt{target\_word = three} \\
\texttt{guess = crane} \\
\texttt{guess\_feedback = c<red> r<yellow> a<red> n<red> e<green>} \\
            }
        }
    }
    & & \\ \\

    \theutterance \stepcounter{utterance}  
    & & & \multicolumn{4}{p{0.6\linewidth}}{
        \cellcolor[rgb]{0.9,0.9,0.9}{
            \makecell[{{p{\linewidth}}}]{
                \texttt{\tiny{[P1$\langle$GM]}}
                \texttt{guess\_feedback: c<red> r<yellow> a<red> n<red> e<green>} \\
\\ 
\texttt{Provide your response only in this format.} \\
\texttt{explanation: details} \\
\texttt{guess: word} \\
            }
        }
    }
     \\ \\

    \theutterance \stepcounter{utterance}  
    & \multicolumn{4}{p{0.6\linewidth}}{
        \cellcolor[rgb]{0.561,0.737,0.561}{
            \makecell[{{p{\linewidth}}}]{
                \texttt{\tiny{[P1$\rangle$GM]}}
                \texttt{explanation: The feedback shows 'r' is in the word but in the wrong position (yellow), while 'a' and 'c' are incorrect (red). 'n' is also incorrect (red), but 'e' is correct and in the right position (green). Since 'r' is misplaced, I will try moving it to a different spot, perhaps after the 'e', and keep 'e' fixed. I'll also replace the incorrect 'n' with another common consonant like 't'.} \\
\texttt{guess: erate} \\
            }
        }
    }
    & & \\ \\

    \theutterance \stepcounter{utterance}  
    & & & \multicolumn{2}{p{0.3\linewidth}}{
        \cellcolor[rgb]{0.9,0.9,0.9}{
            \makecell[{{p{\linewidth}}}]{
                \texttt{\tiny{[GM$|$GM]}}
                \texttt{Error: The guessed word is not a valid word for this game.} \\
            }
        }
    }
    & & \\ \\

    \theutterance \stepcounter{utterance}  
    & & & \multicolumn{4}{p{0.6\linewidth}}{
        \cellcolor[rgb]{0.9,0.9,0.9}{
            \makecell[{{p{\linewidth}}}]{
                \texttt{\tiny{[P1$\langle$GM]}}
                \texttt{Your guess is not a valid word for this game. Please try again.} \\
\\ 
\texttt{Provide your response only in this format.} \\
\texttt{explanation: details} \\
\texttt{guess: word} \\
            }
        }
    }
     \\ \\

    \theutterance \stepcounter{utterance}  
    & \multicolumn{4}{p{0.6\linewidth}}{
        \cellcolor[rgb]{0.561,0.737,0.561}{
            \makecell[{{p{\linewidth}}}]{
                \texttt{\tiny{[P1$\rangle$GM]}}
                \texttt{explanation: The previous guess "erate" was invalid. Based on the feedback from "crane" (r is yellow, e is green), I need a valid word with 'e' at the end and 'r' somewhere else. I will try "order", which is a common word, keeps 'e' at the end, and places 'r' in a different position to test its location further.} \\
\texttt{guess: order} \\
            }
        }
    }
    & & \\ \\

    \theutterance \stepcounter{utterance}  
    & & & \multicolumn{2}{p{0.3\linewidth}}{
        \cellcolor[rgb]{0.9,0.9,0.9}{
            \makecell[{{p{\linewidth}}}]{
                \texttt{\tiny{[GM$|$GM]}}
                \texttt{attempts = 2} \\
\texttt{target\_word = three} \\
\texttt{guess = order} \\
\texttt{guess\_feedback = o<red> r<yellow> d<red> e<green> r<red>} \\
            }
        }
    }
    & & \\ \\

    \theutterance \stepcounter{utterance}  
    & & & \multicolumn{4}{p{0.6\linewidth}}{
        \cellcolor[rgb]{0.9,0.9,0.9}{
            \makecell[{{p{\linewidth}}}]{
                \texttt{\tiny{[P1$\langle$GM]}}
                \texttt{guess\_feedback: o<red> r<yellow> d<red> e<green> r<red>} \\
\\ 
\texttt{Provide your response only in this format.} \\
\texttt{explanation: details} \\
\texttt{guess: word} \\
            }
        }
    }
     \\ \\

    \theutterance \stepcounter{utterance}  
    & \multicolumn{4}{p{0.6\linewidth}}{
        \cellcolor[rgb]{0.561,0.737,0.561}{
            \makecell[{{p{\linewidth}}}]{
                \texttt{\tiny{[P1$\rangle$GM]}}
                \texttt{explanation: The feedback from "order" shows 'r' is definitely in the word but not in the second or fourth position (both red/yellow confusion resolved: first 'r' was yellow in 'crane', now both 'r's in 'order' are red except the one that was yellow before? Wait, let's re{-}evaluate. In 'crane': c(red), r(yellow), a(red), n(red), e(green). So 'r' is in the word, not at index 1. In 'order': o(red), r(yellow), d(red), e(green), r(red). This means 'r' is in the word, not at index 1 (the first 'r' in order) and not at index 4 (the second 'r' in order). Wait, the feedback says r<yellow> for the second letter and r<red> for the fifth letter. This implies the 'r' at index 1 is correct letter but wrong spot, and the 'r' at index 4 is incorrect? No, usually duplicate letters in guess are handled carefully. If the target has one 'r', and I guess two, one might be marked yellow and one red if the yellow one is the "correct" match for the misplaced one. But here, 'r' was yellow at index 1 in 'crane' and yellow at index 1 in 'order'. This suggests 'r' is NOT at index 1. Let's re{-}read the feedback carefully: "o<red> r<yellow> d<red> e<green> r<red>". This is tricky. If 'r' is yellow at index 1, it means 'r' is in the word but not at index 1. If 'r' is red at index 4, it means that specific 'r' is not the match. So 'r' is in the word, not at index 1, not at index 4. We know 'e' is at index 4 (green). We need a word with 'e' at the end, 'r' somewhere else (not 1, not 4), and no 'o', 'd', 'a', 'c', 'n'. Common words ending in 'er' or 're'? 're' is not 5 letters. Words like 'baker', 'maker', 'saver', 'racer'? 'racer' has 'r' at start and end. 'baker': b, a, k, e, r. 'a' was red in crane. 'k' is unknown. 'saver': s, a, v, e, r. 'a' is red. 'rider': r, i, d, e, r. 'd' is red. 'rover': r, o, v, e, r. '} \\
            }
        }
    }
    & & \\ \\

    \theutterance \stepcounter{utterance}  
    & & & \multicolumn{2}{p{0.3\linewidth}}{
        \cellcolor[rgb]{0.9,0.9,0.9}{
            \makecell[{{p{\linewidth}}}]{
                \texttt{\tiny{[GM$|$GM]}}
                \texttt{Error: The response should contain the 'guess:' keyword exactly once.} \\
            }
        }
    }
    & & \\ \\

    \theutterance \stepcounter{utterance}  
    & & & \multicolumn{2}{p{0.3\linewidth}}{
        \cellcolor[rgb]{0.9,0.9,0.9}{
            \makecell[{{p{\linewidth}}}]{
                \texttt{\tiny{[GM$|$GM]}}
                \texttt{game\_result = ABORT} \\
            }
        }
    }
    & & \\ \\

    \theutterance \stepcounter{utterance}  
    & & \multicolumn{4}{p{0.6\linewidth}}{
        \cellcolor[rgb]{0.95,0.95,0.95}{
            \makecell[{{p{\linewidth}}}]{
                \texttt{\tiny{[P1$\leftrightarrow$P]}}
                \texttt{Post{-}episode Reflexion (not sent to GM):} \\
\texttt{New plan: I failed because my explanation became a long, rambling monologue that violated the "concise" rule and likely exceeded the token limit or formatting constraints, causing the Game Master to reject the message. Additionally, I got stuck in a loop of over{-}analyzing the duplicate 'r' feedback without proposing a concrete, valid word immediately. Next time, I will strictly adhere to the format: provide a single, concise sentence in the "explanation:" tag that directly links the feedback to the next move, followed immediately by a valid five{-}letter word in the "guess:" tag. I will avoid internal monologues or re{-}evaluating the logic in the explanation text. Based on the feedback (e at end, r not at index 1 or 4, no o/d/a/c/n), I will guess "baker" (b{-}a{-}k{-}e{-}r) or "saver" (s{-}a{-}v{-}e{-}r) but since 'a' is red, I will try "baker" only if I'm sure 'a' is red, wait, 'a' was red in 'crane'. So 'baker' has 'a'. Bad. How about "baker" is out. "baker" has 'a'. "saver" has 'a'. "rider" has 'd'. "rover" has 'o'. "rebel"? r{-}e{-}b{-}e{-}l (two e's). "reign"? no. "renew"? n is red. "revel"? v is unknown. "revel" {-}> r(yellow in crane? no, r was yellow in crane at index 1, meaning r is in word but not index 1. In 'revel', r is at index 0. That works. e is at index 1 (green requires index 4). So 'revel' fails. I need 'e' at the end. Words ending in 'er' where 'r' is not at 1 or 4. 'r' is at 4 in 'er'. Wait, index 4 is the last letter. If the word ends in 'er', the 'r' is at index 4. But feedback said 'r' at index 4 in 'order' was red. So the word does NOT end in 'r'. But 'e' is green at index 4. So the word ends in 'e'. And 'r' is somewhere else (not 1, not 4). So pattern: ? ? ? ? e. With 'r' in positions 0, 2, or 3. And no o, d, a, c, n. Try "breeze"? b{-}r{-}e{-}e{-}z{-}e (6 letters). "brine"? n is red. "brake"? a is red. "broke"? o is red. "breech"? c is red.} \\
            }
        }
    }
    & \\ \\

\end{supertabular}
}

\end{document}
